\section{Discussion}
\label{sec:discussion}

\subsection{What the study shows}

The research question had two halves, and they do not point the same way.

On the first, the pre-registered comparison is unambiguous within its scope. The full mechanism beat enforced
record-and-temporal labels by $17.1$\pp{} of task success, advisory temporal labels by $14.5$\pp, an LLM arbiter
under the same gated policy by $15.9$\pp, and unlabelled memory by $15.3$\pp, with every interval clear of zero
after Holm correction, in the same direction for all six models, and stable across three item-set definitions.
Harm fell with it: 3.1\% of \cond{EBG} runs were harmful against 22.6\% under \cond{RT}, and arbitrary picks fell
from 252 to 13.

On the second, the answer is yes --- and it cuts against our own design. Enforcing every arbitration rather than
only strong ones raised false resolution from 1.9\% to 12.9\%, restored arbitrary picks from 13 to 166, and more
than doubled the harmful rate. It also scored $9.7$\pp{} \emph{higher} on task success. Gating buys safety with
accuracy; it does not dominate.

\subsection{Where the effect actually comes from}

The ablations locate the benefit more precisely than the headline contrasts do, and they relocate it. Most of it
is in the labels: evidence-basis statuses with no enforcement at all are worth $+11.8$\pp{} over metadata labels
and $+8.5$\pp{} over advisory temporal labels, and cut harm from about 7.5\% to 2.7\%. Gating adds a further
$+6.0$\pp. Safe submission adds nothing measurable to success while halving format errors, and
search-before-submit adds about $1.3$\pp.

The \cond{LG} $-$ \cond{LA} contrast says why that ordering matters. The same gating machinery driven by an LLM
arbiter produced no success gain and $8.1$\pp{} \emph{more} harm, with 37.0\% false resolution. Enforcement
inherits the reliability of whatever supplies its labels: enforcing a record-metadata scanner's labels tripled harm
relative to plain memory, enforcing a language model's tripled it too, and enforcing a deterministic,
evidence-grounded arbiter's strong resolutions did not. The contribution is therefore best stated as \emph{the
evidence basis is what makes enforcement safe to apply and tells you where to stop applying it} --- not that
gating, or enforcement, is good in itself.

\subsection{The label-harm finding replicates}

The study began from a development observation: metadata-derived reliability labels made small agents worse than
no labels, because the agents followed the label instead of the text. On sealed confirmatory data this replicates
without modification --- \cond{G} 43.5\% against \cond{B} 46.0\%, with more harm and more false conflicts --- and
enforcing those labels made it worse again. A reliability signal the evidence does not support is not merely
useless to a small agent; it is actively misleading, and enforcement amplifies it. This is consistent with prior
mechanistic evidence that models can bind a stale value while still encoding the update, and that compliance with
an advisory correction depends on the model family \citep{guo2026cicm}.

\subsection{Where it does not generalise}

Model family matters more than model size for \emph{how} the mechanism helps. For the three Qwen models the story
is the one the hypotheses describe: harmful answers become correct ones, with harm differences of $-24.5$ to
$-27.7$\pp. For Llama-3.1-8B it is not: that model abstains on 44\% of \cond{EBG} runs, fails to format an answer
on 11\%, is the least harmful model in the study under every condition, and gains its $+6.8$\pp{} by converting
abstentions into answers. Its harm difference is $-0.3$\pp{} because there was almost no harm to remove. The
pre-registration anticipated this risk; eval3 gives a small positive effect whose interval excludes zero, which is
better than the development data suggested but is not the same effect as the one measured on Qwen. We do not claim
universal superiority across families.

Nor is the absolute picture good. The best condition answers 61.3\% of tasks correctly and 46.0\% of those with a
definite answer. These are small models on synthetic, template-regular, English, ISO-dated records in one harness,
and on a benchmark that by construction contains the lexically signalled cases this arbiter is built to read.

\subsection{Three kinds of evidence, kept apart}

Three kinds of claim appear in this paper and should not be read alike. The pilot observations and the development
split results are \textbf{development evidence}: they motivated the design and the power analysis and are not
tests of it. The four Holm-corrected contrasts and S4 are \textbf{confirmatory}: pre-registered, frozen before the
data existed, reported here with whatever sign they came out. Everything in \Cref{sec:ablations}, the three
item-set sensitivity analyses and the two platform analyses are \textbf{exploratory}; the last two are
additionally post-hoc. We draw no confirmatory conclusion from any of them --- including the
\cond{EBG} $-$ \cond{EBG-all} trade that we consider the most interesting result in the paper.

\subsection{What would change our mind}

Four results would undermine the contribution as stated. A replication in which evidence-basis labels without
enforcement did not beat metadata labels would remove the part of the effect we consider best supported. A
benchmark that prices wrong confident answers against missed correct ones, on which ungated enforcement came out
ahead overall, would make gating the wrong default rather than a defensible one. A model family on which
\cond{EBG} lowered success would bound the claim further than Llama-3.1-8B already does. And a transcript study
showing that \cond{EBG}'s remaining 20.2\% false conflicts are items the arbiter \emph{could} have resolved from
the text would move its known limitations from a documented residue to a design flaw.
